\newcommand{\F}{\mathbb{F}}
\newcommand{\C}{\mathbb{C}}
\newcommand{\an}{\mathscr{A}} 
\newcommand{\Q}{\mathbb{Q}}
\renewcommand{\hom}{\operatorname{Hom}_\F}
\newcommand{\ehom}{\operatorname{End}_\F}
\colorlet{gris}{black!60}

\-\vspace{-0.3cm}
  {\hrule \vspace{-0.15cm}
  \centering \textbf{P1} \vspace{0.1cm}
  \hrule}

3. Determine se cada uma das afirmações é abaixo é verdadeira ou falsa. \\
(a) Se existem \(u,v \in V \) l.i \(: m_u = m_v= t^3\), então \(t^4 \ |\ c_{_T}\). \\ 
\textcolor{gris}{FALSO. 
} \\ 
(b) Se \(V_1, \ldots, V_m\) são os subespaços \(T\)-primários de \(V\) e \(W\leq V\) é \(T\)-inv, então \(W = \bigoplus^m (W\cap V_j)\). \\ 
\textcolor{gris}{VERDADEIRO. 
} \\ 
(c) Se \(T \in \ehom(V)\) é injetor, então \(T\) tambem é sobrejetor. \\ 
\textcolor{gris}{FALSO. 
}\\ 
4. De um exemplo de operador linear num \(\F\)-espaço \(V\) com \(3\) fatores invariantes pra o qual existam vetores \(v_1, v_2, v_3 \in V\) satisfazendo: 
\begin{enumerate}[label = \roman*.]
\item \(V = C_T(v_1) \oplus C_T(v_2) \oplus C_T(v_3)\).
\item \(m_{v_1} \) não divide \(m_{v_2}\) nem \(m_{v_3}\). 
\item \(m_{v_2}\) não divide \(m_{v_1}\). 
\end{enumerate}
